\chapter{Mapa de dependencias entre resultados, tareas e investigación}

\section{Qué significa una flecha en este mapa}

\begin{context}
Una flecha continua indica que el resultado de origen se usa en la demostración o en la
construcción del resultado de destino. Una flecha discontinua indica una aplicación o comparación
conceptual: no afirma que el resultado de destino se deduzca del primero. Las hipótesis escritas
junto a cada ruta son parte de la conexión.
\end{context}

\section{Rutas que intervienen en demostraciones}

\begin{dependencias}
\textbf{Continuidad e inversos.}
\[
\boxed{\text{Serie de Neumann (Parte 1)}}
\longrightarrow
\boxed{\text{T1.1(b): inversión continua}}
\longrightarrow
\boxed{\text{Prop. 2.17: permanencia espectral}}.
\]
En la última flecha se usa que $(b+i/n)^{-1}\to b^{-1}$ en la álgebra ambiente y que la
$C^*$-subálgebra es cerrada. T1.1(a) aporta, por otra vía, continuidad del producto para T1.3,
T1.7 y T2.8.
\end{dependencias}

\begin{dependencias}
\textbf{Reducción conmutativa.}
\[
\begin{gathered}
\boxed{\text{T1.2: inversos que conmutan}}
\quad+\quad
\boxed{\text{T1.3: álgebra generada}}\\[2mm]
\Downarrow\\[2mm]
\boxed{\text{T1.9: subálgebra inversamente cerrada}}\\[2mm]
\Downarrow\\[2mm]
\boxed{\text{Prop. 2.13: espectro real}}.
\end{gathered}
\]
La última prueba añade la identidad $C^*$, la realidad de los caracteres en autoadjuntos y la
descripción del espectro mediante caracteres. T1.9 aporta la subálgebra conmutativa que conserva
el espectro; por sí solo no demuestra que el espectro sea real.
\end{dependencias}

\begin{dependencias}
\textbf{De las funciones al cálculo funcional.}
\[
\begin{gathered}
\boxed{\text{T1.3(d): }C^*(a,1)\text{ conmutativa para $a$ normal}}\\[2mm]
\quad+\quad
\boxed{\text{Prop. 2.17: permanencia espectral}}
\quad+\quad
\boxed{\text{Teo. 2.20: Gelfand--Naimark}}\\[2mm]
\Downarrow\\[2mm]
\boxed{\text{Teo. 2.22: cálculo funcional}}.
\end{gathered}
\]
T1.10 calcula el espectro de una función continua y, combinado con el cálculo funcional y la
permanencia espectral, da T2.4. T2.2(a,c) proporciona las direcciones directas de T2.6(c,b);
las recíprocas usan normalidad y cálculo funcional.
\end{dependencias}

\begin{dependencias}
\textbf{Continuidad al variar el operador.}
\[
\begin{gathered}
\boxed{\text{T1.7: límite de normales normal}}
\quad+\quad
\boxed{\text{Stone--Weierstrass}}\\[2mm]
\quad+\quad
\boxed{\text{cálculo funcional isométrico}}\\[2mm]
\Downarrow\\[2mm]
\boxed{\text{T2.8: }f(x_n)\to f(x)}.
\end{gathered}
\]
La convergencia aquí es en norma. T2.9 estudia otra situación: una red creciente y acotada de
autoadjuntos converge fuertemente, es decir, vector por vector en norma de $H$.
\end{dependencias}

\section{Cobertura de los ejercicios}

\small
\begin{longtable}{>{\bfseries}p{1.3cm}p{4.3cm}p{8.2cm}}
\toprule
Tarea & Resultado demostrado & Uso exacto o alcance de la conexión\\
\midrule
\endhead
T1.1 & Producto e inversión continuos & El producto permite T1.3 y T1.7; la inversión se usa en la Prop. 2.17 y el producto de nuevo en T2.8.\\
T1.2 & Conmutación con inversos y adjuntos & Se usa al añadir inversos en la construcción de T1.9.\\
T1.3 & Álgebra generada y conmutación & Para un normal $a$, construye $C^*(a,1)$ conmutativa, usada en las Props. 2.18 y 2.22.\\
T1.4 & Cociente por ideal cerrado & Hace de $\mathcal B(X)/\K(X)$ una álgebra de Banach; la equivalencia Fredholm--invertibilidad es el teorema externo de Atkinson.\\
T1.5 & $C_0(X)$ cerrado en $C_b(X)$ & Si $X$ no es compacto, es no unital; la clasificación general de $C_0(X)$ requiere la versión no unital del teorema de la Sección 2.\\
T1.6 & $\spec(x^*)=\overline{\spec(x)}$ & Simetría válida en álgebras involutivas; la realidad del espectro autoadjunto exige además el argumento $C^*$ de la Prop. 2.13.\\
T1.7 & Normalidad cerrada & Primera hipótesis que T2.8 necesita para definir $f(x)$ en el límite.\\
T1.8 & Inclusión espectral por homomorfismos & Aplicada a $a^*a$, entra en la prueba de contractividad de $*$-homomorfismos (Teo. 2.8).\\
T1.9 & Subálgebra conmutativa inversamente cerrada & Se usa expresamente en la prueba de la Prop. 2.13 para preservar el espectro de un autoadjunto.\\
T1.10 & $\spec(f)=f(X)$ para $X$ compacto & Calcula el espectro de $f$ en la prueba de T2.4; ilustra el modelo $C(X)$ de Gelfand--Naimark.\\
T1.11 & $\mathcal B(H)$ y shift unilateral & El inciso (a) da el modelo concreto de T2.1; (c,d) separan espectro puntual aproximado y espectro, y T2.10 completa el cálculo de $\spec(S)$.\\
T1.12 & Álgebra del disco & Muestra que una inclusión unital y cerrada de álgebras de Banach puede cambiar el espectro; la inclusión no es de $C^*$-subálgebras.\\
\midrule
T2.1 & $C_b(X)\hookrightarrow\mathcal B(\ell^2(X))$ & Verifica directamente un $*$-homomorfismo isométrico; ejemplo concreto del Teo. 2.27.\\
T2.2 & Idempotentes, productos, unitarios & (a,c) dan direcciones directas de la Prop. 2.28; (b) estima radios de productos conmutativos y no se necesita para esa proposición.\\
T2.3 & Producto acotado y suma $c_0$ & Construye una $C^*$-álgebra y su ideal cerrado; T1.4 permite considerar el cociente como álgebra de Banach.\\
T2.4 & Mapeo espectral continuo & Depende del Teo. 2.22, la Prop. 2.17 y T1.10; se aplica al espectro de funciones de normales.\\
T2.5 & $X\cong\Omega(C(X))$ & Da la unicidad del espacio compacto en el Cor. 2.21.\\
T2.6 & Autoadjuntos, unitarios y proyecciones normales & Reconstruye la Prop. 2.28 mediante cálculo funcional y muestra con contraejemplos la necesidad de normalidad.\\
T2.7 & Factor unitario de un invertible & Usa la raíz positiva de la Prop. 2.30 y la invertibilidad de $a^*a$ para obtener $a=u|a|$.\\
T2.8 & Continuidad de $f(x_n)$ & Combina T1.7, Stone--Weierstrass y cálculo funcional; la convergencia es en norma.\\
T2.9 & Convergencia monótona de autoadjuntos & Usa orden positivo, polarización y Riesz; concluye convergencia fuerte, sin afirmar convergencia en norma.\\
T2.10 & Espectro del shift unilateral & Amplía T1.11(d) a $\spec(S)=\overline{\mathbb D}$; $S$ no es normal.\\
\bottomrule
\end{longtable}
\normalsize

\section{Dos puentes controlados hacia Haseman}

\begin{center}
\begin{tikzpicture}[node distance=8mm]
\node[taskbox,text width=2.6cm] (q) {T1.4\\cociente por ideal};
\node[theorybox,text width=2.7cm,right=of q] (calkin) {$\K(X)$ cerrado\\$\mathcal B(X)/\K(X)$};
\node[researchbox,text width=2.6cm,right=of calkin] (class) {Clase de $A$\\$A+\K(X)$};
\node[flowend,text width=2.7cm,right=of class] (fred) {Atkinson\\invertible $\Leftrightarrow$ Fredholm};
\draw[arr] (q)--(calkin);
\draw[arr] (calkin)--(class);
\draw[arr] (class)--(fred);
\end{tikzpicture}
\end{center}

\noindent Para $X$ de Banach e infinito-dimensional, la primera flecha usa que $\K(X)$ es un ideal cerrado. La última es el teorema de
Atkinson, no una consecuencia de T1.4. En los operadores de la línea Haseman, esta formulación
permite plantear el problema de Fredholm en un cociente; véase
\cite[p.~2]{Karlovich2025}.
Como comprobación concreta, en $X=\ell^2$ tome $P$ como proyección
sobre el primer vector canónico $e_1$ y ponga $A=I_X-P$.
Entonces $A e_1=0$, pero
$A+\K(X)=I_X+\K(X)$ porque $P$ es compacto. El cálculo muestra qué
información conserva el cociente y por qué T1.4 prepara la formulación
de Atkinson.

\begin{center}
\begin{tikzpicture}[node distance=9mm and 8mm]
\node[theorybox,text width=3.0cm] (model) {Modelo $L^2(\Gamma_0)$,\\$\Gamma_0=\mathbb R$ o $\mathbb T$\\$S_{\Gamma_0}^2=I$, $S_{\Gamma_0}^*=S_{\Gamma_0}$};
\node[taskbox,text width=3.0cm,right=of model] (orth) {$P_{\Gamma_0}^\pm=(I\pm S_{\Gamma_0})/2$\\proyecciones ortogonales};
\draw[arr] (model)--(orth);
\node[researchbox,text width=3.0cm,below=of model] (plem) {Estrella abierta de Karlovich\\Sokhotski--Plemelj};
\node[researchbox,text width=3.0cm,right=of plem] (trace) {$P_\pm=(I\pm S_\Gamma)/2$\\componentes de trazas};
\node[flowend,text width=3.0cm,right=of trace] (has) {$T=V_\alpha P_+ + GP_-$\\ecuación de Haseman};
\draw[arr] (plem)--(trace);
\draw[arr] (trace)--(has);
\draw[mergearr,dash pattern=on 4pt off 2pt] (orth.south)--(trace.north);
\end{tikzpicture}
\end{center}

\noindent La flecha discontinua sólo compara la forma de las fórmulas.
En el modelo de Bastos--Fernandes--Karlovich \cite[§2, p.~22]{BFK2006},
$P_{\Gamma_0}^\pm$ son proyecciones ortogonales. En la estrella abierta de
Karlovich 2017, la reducción de frontera a $T$ procede de
Sokhotski--Plemelj \cite[pp.~463--464]{Karlovich2017}; la identidad
$S_\Gamma^2=I$ no vale en general allí, según el teorema 5.3 de esa fuente.
Por tanto, los operadores $P_\pm$ de esa ecuación no se llaman aquí
idempotentes ni proyecciones.
En el caso de comprobación $\alpha=\mathrm{id}$ y $G=1$, la fórmula de
Haseman da $T=P_++P_-=I$ por cancelación de los dos términos $S_\Gamma$.
Este cálculo usa la fórmula de Plemelj, sin requerir $S_\Gamma^2=I$.

\begin{reflection}
Para comprobar una flecha, localiza la línea de la demostración que usa su origen. Por ejemplo,
identifica dónde entra T1.9 en la realidad del espectro y dónde entra T1.10 en T2.4. Después
verifica en el modelo $L^2(\Gamma_0)$ cómo $S_{\Gamma_0}^2=I$ produce
idempotentes y cómo $S_{\Gamma_0}^*=S_{\Gamma_0}$ los hace proyecciones
ortogonales. Contrasta esas hipótesis con las de la estrella abierta.
\end{reflection}
